% SPDX-License-Identifier: GPL-3.0-or-later
% Copyright (C) 2026 Alessandro Greco
\chapter{Руководство по спецификациям torspec}\label[appendice]{cap:C-specifiche}
\aggiornatoal{23 сентября 2026 г.}{коммит torspec \texttt{8e1f245}}

Спецификации Tor Project являются основным источником почти для каждой главы этой книги. Они исчерпывающи, но создавались не как учебник. Более двадцати лет их поддерживали как отдельные текстовые файлы, и лишь сравнительно недавно объединили на одном сайте. Даже собственные соглашения документации признают, что работа над единообразием ещё продолжается. Это приложение служит картой для читателей, которые после книги хотят перейти к непосредственному изучению спецификаций.

% ---------------------------------------------------------------------------
\section{Где их найти}\label{sec:C-dove}

Спецификации опубликованы по адресу \url{https://spec.torproject.org/} и поддерживаются вместе с предложениями по развитию протокола в GitLab-репозитории \texttt{torspec} Tor Project. В репозитории спецификации находятся в Markdown-файлах внутри \file{spec/}: небольшим темам соответствуют отдельные файлы, более крупным документам — каталоги. Файл \file{permalinks.md} перечисляет стабильные адреса документов, например \url{https://spec.torproject.org/tor-spec}; при цитировании следует использовать именно их. Во всей книге документы цитируются как \spec{document}{section}, причём название раздела сохраняется таким, каким оно приведено на сайте, а предложения — как \enquote{Предложение~NNN}.

% ---------------------------------------------------------------------------
\section{Соглашения}\label{sec:C-convenzioni}

Вводный документ с соглашениями устанавливает правила, действующие для спецификаций, если явно не указано иное:
\begin{itemize}
  \item \texttt{MUST}, \texttt{SHOULD}, \texttt{MAY} и их отрицательные формы имеют значения RFC~2119: соответственно обязательное требование, сильная рекомендация или разрешённая возможность.
  \item Длины измеряются в восьмибитных байтах, а двоичные префиксы вроде \enquote{kibibyte} означают кратные 1024.
  \item Многобайтовые целые записываются в big-endian, или \enquote{сетевом}, порядке байтов: десятичное число 4660 в двух байтах записывается как \texttt{12 34}.
\end{itemize}
Два важных исключения уже встречались в книге: ключи и точки Curve25519 и Ed25519 используют little-endian (\cref{cap:B-crittografia}), а некоторые документы системы каталогов имеют текстовый формат с двоичными данными, закодированными в base64.

% ---------------------------------------------------------------------------
\section{Карта документов}\label{sec:C-mappa}

\begin{center}\small
\begin{tabularx}{\textwidth}{@{}l>{\raggedright\arraybackslash}Xl@{}}
\toprule
Документ & Содержимое & Главы \\
\midrule
\texttt{tor-spec} & Основной протокол: каналы, ячейки, цепи, рукопожатия, послойное шифрование на ретрансляторах, потоки, управление потоком и перегрузкой & \ref{cap:07-canali}--\ref{cap:11-stream} \\
\texttt{cert-spec} & Сертификаты Ed25519 в Tor & \ref{cap:07-canali} \\
\texttt{dir-spec} & Центры каталогов, голосования, консенсус, дескрипторы, микродескрипторы и документы extra-info & \ref{cap:12-directory}, \ref{cap:24-metrics} \\
\texttt{path-spec} & Выбор маршрута и ограничения & \ref{cap:13-percorso} \\
\texttt{guard-spec} & Алгоритм выбора входных узлов (guard) & \ref{cap:13-percorso} \\
\texttt{vanguards-spec} & Полноценные и облегчённые защитные промежуточные узлы (vanguards) & \ref{cap:13-percorso} \\
\texttt{rend-spec} & Onion-сервисы v3: ключи, дескрипторы, точки представления и встречи & \ref{cap:14-onion-services}, \ref{cap:21-onion-pratica} \\
\texttt{hspow-spec} & Доказательство выполнения работы (PoW) для onion-сервисов & \ref{cap:15-dos} \\
\texttt{dos-spec} & Защита от отказа в обслуживании (DoS) & \ref{cap:15-dos} \\
\texttt{pt-spec} & Интерфейс между tor и подключаемыми транспортами & \ref{cap:16-censura} \\
\texttt{padding-spec} & Заполнение каналов и цепей & \ref{cap:07-canali}, \ref{cap:26-attacchi} \\
\texttt{control-spec} & Протокол ControlPort & \ref{cap:18-stem} \\
\texttt{socks-extensions} & Специфичные для Tor расширения SOCKS & \ref{cap:11-stream}, \ref{cap:28-opsec} \\
\texttt{param-spec} & Сетевые параметры консенсуса и их значения по умолчанию & \ref{cap:11-stream}, \ref{cap:15-dos} \\
\texttt{address-spec} & Специальные имена, например адреса \texttt{.onion} & \ref{cap:14-onion-services} \\
\texttt{srv-spec} & Совместное случайное значение центров каталогов & \ref{cap:14-onion-services} \\
\bottomrule
\end{tabularx}
\end{center}

% ---------------------------------------------------------------------------
\section{Как читать раздел спецификации}\label{sec:C-lettura}

Несколько рекомендаций, выработанных при написании книги:
\begin{enumerate}
  \item \textbf{Начинайте с формата.} Разделы о бинарных сообщениях приводят имена и размеры полей в таблицах. Прежде чем читать правила, восстановите формат на бумаге, как мы делали в схемах форматов книги.
  \item \textbf{Отличайте актуальный материал от устаревшего.} Многие документы сохраняют исторические функции, например TAP, версии протокола связи 1--3 и onion-сервисы v2. Устаревший материал часто сгруппирован в разделах или файлах вроде \file{obsolete-channels.md} либо \file{historical/}, но не всегда. При сомнении исследуйте код и проверьте, существует ли соответствующая ветвь реализации.
  \item \textbf{Сравнивайте с реализациями.} Спецификации выражают намерение; кодовые базы tor и Arti показывают фактическое поведение. Во время исследований для книги несколько расхождений оказались особенно информативными: отсутствие заполнения канала в одноузловых цепях (\cref{cap:07-canali}) и наблюдаемая на выходном ретрансляторе причина закрытия потока (\cref{sec:lab-relay}).
  \item \textbf{Проверяйте статус предложений.} Предложение (proposal) может быть реализовано до включения в спецификацию либо всё ещё обсуждаться и не иметь реализации (\cref{sec:fut-processo}).
  \item \textbf{Проверяйте экспериментально.} Приватная сеть вроде лаборатории книги позволяет проверить, ведёт ли себя поле так, как описано. На этом методе основаны все лабораторные работы.
\end{enumerate}

% ---------------------------------------------------------------------------
\section{Arti и спецификации}\label{sec:C-arti}

Код Arti часто ссылается на спецификации в комментариях и документации и отмечает неоднозначные требования или различия с tor. Таблица поддерживаемых подпротоколов (\cref{sec:arti-supporto}) показывает, какие функции поддерживают обе реализации и, следовательно, какие части спецификаций описывают \enquote{живой} протокол.
